\subsection{Strongly relatively prime polynomials}

Let R be a commutative ring. Two elements x, y of R are called strongly relatively prime if the principal ideals Rx and Ry are relatively prime, in other words (Chapter II, § 1, no. 2) if \(\mathbf{R}x + \mathbf{R}y = \mathbf{R}\) ; it amounts to the same to say that there exist two elements \(a, b\) of \(\mathbf{R}\) such that \(ax + by = 1\) .

LEMMA 1 (``Euclid's Lemma''). Let \(x, y\) be two strongly relatively prime elements of \(\mathbf{R}\) ; if \(z \in \mathbf{R}\) is such that \(x\) divides \(yz\) , then \(x\) divides \(z\) .

If \(1 = ax + by\) , then \(z = x(az) + (yz)b\) .

If \(x\) and \(y\) are strongly relatively prime in \(\mathbf{R}\) , then

\[
\mathrm{R} x y = (\mathrm{R} x) \cap (\mathrm{R} y)
\]

(Chapter II, § 1, no. 2, Proposition 5); if R is an integral domain, two strongly relatively prime elements then have an lcm equal to their product (Algebra, Chapter VI, § 1, no. 8) and are therefore relatively prime in the sense of Algebra, Chapter VI, § 1, no. 12. Conversely, if R is a principal ideal domain, two relatively prime elements are also strongly relatively prime, as follows from Bezout's identity (Algebra, Chapter VII, § 1, no. 2, Theorem 1).

For polynomial rings there is the following result:

PROPOSITION 1. Let A be a commutative ring and P and P' two strongly relatively prime polynomials in A{[}X{]}. Suppose that P is monic and of degree s. Then every polynomial T in A{[}X{]} may be written uniquely in the form

\[
\mathrm{T} = \mathrm{PQ} + \mathrm{P} ^ {\prime} \mathrm{Q} ^ {\prime}\tag{1}
\]

where \(\mathbf{Q} \in \mathbf{A}[\mathbf{X}]\) , \(\mathbf{Q}' \in \mathbf{A}[\mathbf{X}]\) and \(\deg(\mathbf{Q}') < s\) .

If further \(\deg (\mathbf{T})\leqslant t\) and \(\deg (\mathbf{P}^{\prime})\leqslant t - s\) , then \(\deg (\mathbf{Q})\leqslant t - s\)

As \(\mathbf{P}\) is monic, \(\mathrm{PR} \neq 0\) for every polynomial \(\mathbf{R} \neq 0\) of \(\mathbf{A}[\mathbf{X}]\) and in this case \(\deg(\mathrm{PR}) = s + \deg(\mathbf{R})\) .

Let \(\mathbf{T}\) be any polynomial in A{[}X{]}. As the ideal generated by P and \(\mathbf{P}'\) is the whole of A{[}X{]}, there exist polynomials \(\mathbf{Q}_1\) and \(\mathbf{Q}_1'\) such that

\[
\mathrm{T} = \mathrm{PQ} _ {1} + \mathrm{P} _ {1} ^ {\prime} \mathrm{Q} _ {1} ^ {\prime};
\]

as P is monic of degree s, Euclidean division (Algebra, Chapter IV, § 1, no. 5) shows that there exist two polynomials Q', Q'\,' such that Q'\_1 = PQ'\,' + Q' where deg(Q') \textless{} s; then we deduce that

\[
\mathrm{T} = \mathrm{PQ} _ {1} + \mathrm{P} ^ {\prime} (\mathrm{PQ} ^ {\prime \prime} + \mathrm{Q} ^ {\prime}) = \mathrm{PQ} + \mathrm{P} ^ {\prime} \mathrm{Q} ^ {\prime}
\]

where \(Q = Q_{1} + P'Q''\) . To show the uniqueness of formula (1), it is sufficient to prove that the relations

\[
0 = \mathrm{PQ} + \mathrm{P} ^ {\prime} \mathrm{Q} ^ {\prime}, \quad \deg (\mathrm{Q} ^ {\prime}) <   s\tag{2}
\]

imply \(Q = Q' = 0\) . Now, if (2) holds, \(P\) divides \(-PQ = P'Q'\) and, as \(P\) and

\(P'\) are strongly relatively prime, \(P\) divides \(Q'\) by Lemma 1; if \(Q' \neq 0\) , there would exist a polynomial \(S \neq 0\) such that \(Q' = PS\) , whence

\[
\deg (\mathbf {Q} ^ {\prime}) = s + \deg (\mathbf {S}) \geqslant s,
\]

which is a contradiction. We conclude that \(Q' = 0\) , whence PQ = 0 and finally Q = 0 by the remark at the beginning.

Finally, suppose that \(\deg(T) \leqslant t\) and \(\deg(P') \leqslant t - s\) ; with the polynomial \(T\) in the form (1),

\[
\deg (\mathrm{P} ^ {\prime} \mathrm{Q} ^ {\prime}) \leqslant \deg (\mathrm{P} ^ {\prime}) + \deg (\mathrm{Q} ^ {\prime}) <   s + \deg (\mathrm{P} ^ {\prime}) \leqslant t
\]

and therefore

\[
s + \deg (Q) = \deg (P Q) = \deg (T - P ^ {\prime} Q ^ {\prime}) \leqslant t
\]

whence \(\deg (\mathbb{Q})\leqslant t - s.\)

Example. For a polynomial \(P \in A[X]\) to be strongly relatively prime to X - a (where \(a \in A\) ), it is necessary and sufficient that \(P(a)\) be invertible in A. For if P and X - a are strongly relatively prime, it follows from Proposition 1 that there exist \(c \in A\) and a polynomial \(Q \in A[X]\) such that \(cP + (X - a)Q = 1\) , whence \(cP(a) = 1\) and \(P(a)\) is invertible. Conversely, by Euclidean division

\[
\mathbf {P} = (\mathbf {X} - a) \mathbf {R} + \mathbf {P} (a)
\]

and, if \(\mathrm{P}(a)=b^{-1}\) , where \(b\in A\) , we deduce that \(1=b\mathrm{P}-b(\mathrm{X}-a)\mathrm{R}\) , which shows that P and X-a are strongly relatively prime.

Let A and B be two commutative rings and \(f\colon\mathbf{A}\to\mathbf{B}\) a ring homomorphism. If \(\mathbf{P}=\sum_{i\geqslant0}a_{i}\mathbf{X}^{i}\) is a formal power series in \(\mathbf{A}[[\mathbf{X}]]\) , let \(\bar{f}(\mathbf{P})\) denote the formal power series \(\sum_{i\geqslant0}f(a_{i})\mathbf{X}^{i}\) in \(\mathbf{B}[[\mathbf{X}]]\) . If \(\mathbf{P}\) is a polynomial, so is \(\bar{f}(\mathbf{P})\) and, if further \(\mathbf{P}\) is monic, then \(\bar{f}(\mathbf{P})\) is monic of the same degree as \(\mathbf{P}\) . Finally, \(\mathbf{P}\mapsto\bar{f}(\mathbf{P})\) is clearly a homomorphism of \(\mathbf{A}[[\mathbf{X}]]\) to \(\mathbf{B}[[\mathbf{X}]]\) which extends \(f\) and maps \(\mathbf{X}\) to \(\mathbf{X}\) . The notation \(\bar{f}\) will be constantly used in this sense in the rest of this paragraph.

PROPOSITION 2. Let A and B be two commutative rings, f a homomorphism from A to B and P, P' two polynomials in A{[}X{]}. If P and P' are strongly relatively prime in A{[}X{]}, then \(\bar{f}(\mathbf{P})\) and \(\bar{f}(\mathbf{P}')\) are strongly relatively prime in B{[}X{]}. The converse is true if \(f\) is surjective, if its kernel is contained in the Jacobson radical of A and if P is monic.

Suppose that P and \(P'\) are strongly relatively prime; then there exist polynomials Q, \(Q'\) in A{[}X{]} such that \(PQ + P'Q' = 1\) ; we deduce that

\[
\bar {f} (\mathbf {P}) \bar {f} (\mathbf {Q}) + \bar {f} (\mathbf {P} ^ {\prime}) \bar {f} (\mathbf {Q} ^ {\prime}) = 1,
\]

whence the first assertion. To show the second, let a denote the kernel of f;

let \(\mathbf{E} = \mathbf{A}[\mathbf{X}]\) and \(\mathbf{F}\) the ideal of \(\mathbf{E}\) generated by \(\mathbf{P}\) and \(\mathbf{P}'\) ; as \(f\) is surjective and \(\bar{f}(\mathbf{P})\) monic, Proposition 1 shows that for every polynomial \(\mathbf{T} \in \mathbf{A}[\mathbf{X}]\) there exist two polynomials \(\mathbf{Q}, \mathbf{Q}'\) in \(\mathbf{A}[\mathbf{X}]\) such that

\[
\bar {f} (\mathrm{T}) = \bar {f} (\mathrm{P}) \bar {f} (\mathrm{Q}) + \bar {f} (\mathrm{P} ^ {\prime}) \bar {f} (\mathrm{Q} ^ {\prime}),
\]

whence the relation \(E = F + aE\) . Now, E/F is a finitely generated A-module, for every polynomial is congruent mod. P to a polynomial of degree \(<\deg(P)\) , P being monic. As \(E/F = a(E/F)\) and a is contained in the Jacobson radical of A, Nakayama's Lemma shows that \(E/F = 0\) (Algebra, Chapter VIII, §6, no. 3, Corollary 2 to Proposition 6), which means that P and \(P'\) are strongly relatively prime.

